%%% Hlavní soubor. Zde se definují základní parametry a odkazuje se na ostatní části. %%%

%% Verze pro jednostranný tisk:
% Okraje: levý 40mm, pravý 25mm, horní a dolní 25mm
% (ale pozor, LaTeX si sám přidává 1in)
\documentclass[12pt,a4paper]{report}
\setlength\textwidth{145mm}
\setlength\textheight{247mm}
\setlength\oddsidemargin{15mm}
\setlength\evensidemargin{15mm}
\setlength\topmargin{0mm}
\setlength\headsep{0mm}
\setlength\headheight{0mm}
% \openright zařídí, aby následující text začínal na pravé straně knihy
\let\openright=\clearpage

%% Pokud tiskneme oboustranně:
% \documentclass[12pt,a4paper,twoside,openright]{report}
% \setlength\textwidth{145mm}
% \setlength\textheight{247mm}
% \setlength\oddsidemargin{15mm}
% \setlength\evensidemargin{0mm}
% \setlength\topmargin{0mm}
% \setlength\headsep{0mm}
% \setlength\headheight{0mm}
% \let\openright=\cleardoublepage

%% Použité kódování znaků: obvykle latin2, cp1250 nebo utf8:
\usepackage[utf8]{inputenc}

%% Ostatní balíčky
\usepackage{graphicx}
\usepackage{amsthm}

%% Balíček hyperref, kterým jdou vyrábět klikací odkazy v PDF,
%% ale hlavně ho používáme k uložení metadat do PDF (včetně obsahu).
%% POZOR, nezapomeňte vyplnit jméno práce a autora.
\usepackage[ps2pdf,unicode]{hyperref}   % Musí být za všemi ostatními balíčky
\hypersetup{pdftitle=Querying NoSQL databases in MPS}
\hypersetup{pdfauthor=Radovan Duga}

%%% Drobné úpravy stylu

% Tato makra přesvědčují mírně ošklivým trikem LaTeX, aby hlavičky kapitol
% sázel příčetněji a nevynechával nad nimi spoustu místa. Směle ignorujte.
\makeatletter
\def\@makechapterhead#1{
  {\parindent \z@ \raggedright \normalfont
   \Huge\bfseries \thechapter. #1
   \par\nobreak
   \vskip 20\p@
}}
\def\@makeschapterhead#1{
  {\parindent \z@ \raggedright \normalfont
   \Huge\bfseries #1
   \par\nobreak
   \vskip 20\p@
}}
\makeatother

% Toto makro definuje kapitolu, která není očíslovaná, ale je uvedena v obsahu.
\def\chapwithtoc#1{
\chapter*{#1}
\addcontentsline{toc}{chapter}{#1}
}

\begin{document}

% Trochu volnější nastavení dělení slov, než je default.
\lefthyphenmin=2
\righthyphenmin=2

%%% Titulní strana práce

\pagestyle{empty}
\begin{center}

\large

Charles University in Prague

\medskip

Faculty of Mathematics and Physics

\vfill

{\bf\Large MASTER THESIS}

\vfill

\centerline{\mbox{\includegraphics[width=60mm]{img/logo.eps}}}

\vfill
\vspace{5mm}

{\LARGE Radovan Duga}

\vspace{15mm}

% Název práce přesně podle zadání
{\LARGE\bfseries Querying NoSQL databases in MPS}

\vfill

% Název katedry nebo ústavu, kde byla práce oficiálně zadána
% (dle Organizační struktury MFF UK)
Department of Distributed and Dependable Systems

\vfill

\begin{tabular}{rl}

Supervisor of the master thesis: & RNDr. Pavel Parízek, Ph.D. \\
\noalign{\vspace{2mm}}
Study programme: & Informatics \\
\noalign{\vspace{2mm}}
Specialization: & ISS \\
\end{tabular}

\vfill

% Zde doplňte rok
Prague 2013

\end{center}


\newpage

%%% Následuje vevázaný list -- kopie podepsaného "Zadání diplomové práce".
%%% Toto zadání NENÍ součástí elektronické verze práce, nescanovat.

%%% Na tomto místě mohou být napsána případná poděkování (vedoucímu práce,
%%% konzultantovi, tomu, kdo zapůjčil software, literaturu apod.)

\openright

\noindent
Dedication.

% TODO: dedication

\newpage

%%% Strana s čestným prohlášením k diplomové práci

\vglue 0pt plus 1fill

\noindent
I declare that I carried out this master thesis independently, and only with the cited
sources, literature and other professional sources.

\medskip\noindent
I understand that my work relates to the rights and obligations under the Act No.
121/2000 Coll., the Copyright Act, as amended, in particular the fact that the Charles
University in Prague has the right to conclude a license agreement on the use of this
work as a school work pursuant to Section 60 paragraph 1 of the Copyright Act.

\vspace{10mm}

\hbox{\hbox to 0.5\hsize{%
In ........ date ............
\hss}\hbox to 0.5\hsize{%
signature of the author
\hss}}

\vspace{20mm}
\newpage

%%% Povinná informační strana diplomové práce

\vbox to 0.5\vsize{
\setlength\parindent{0mm}
\setlength\parskip{5mm}

Název práce:
Dotazování NoSQL databáz v prostředí MPS
% přesně dle zadání

Autor:
Radovan Duga

Katedra:  % Případně Ústav:
Katedra distribuovaných a spolehlivých systémů
% dle Organizační struktury MFF UK

Vedoucí diplomové práce:
RNDr. Pavel Parízek, Ph.D., Katedra distribuovaných a spolehlivých systémů
% dle Organizační struktury MFF UK, případně plný název pracoviště mimo MFF UK

Abstrakt:
% abstrakt v rozsahu 80-200 slov; nejedná se však o opis zadání diplomové práce

Klíčová slova:
NoSQL, MPS, dotaz, Cypher
% 3 až 5 klíčových slov

\vss}\nobreak\vbox to 0.49\vsize{
\setlength\parindent{0mm}
\setlength\parskip{5mm}

Title:
Querying NoSQL databases in MPS
% přesný překlad názvu práce v angličtině

Author:
Radovan Duga

Department:
Department of Distributed and Dependable Systems
% dle Organizační struktury MFF UK v angličtině

Supervisor:
RNDr. Pavel Parízek, Ph.D., Department of Distributed and Dependable Systems
% dle Organizační struktury MFF UK, případně plný název pracoviště
% mimo MFF UK v angličtině

Abstract:
% abstrakt v rozsahu 80-200 slov v angličtině; nejedná se však o překlad
% zadání diplomové práce

Keywords:
NoSQL, MPS, query, Cypher
% 3 až 5 klíčových slov v angličtině

\vss}

\newpage

%%% Strana s automaticky generovaným obsahem diplomové práce. U matematických
%%% prací je přípustné, aby seznam tabulek a zkratek, existují-li, byl umístěn
%%% na začátku práce, místo na jejím konci.

\openright
\pagestyle{plain}
\setcounter{page}{1}
\tableofcontents

%%% Jednotlivé kapitoly práce jsou pro přehlednost uloženy v samostatných souborech
\chapter{Introduction}
% \addcontentsline{toc}{chapter}{Introduction}

- there exist lot of general dsl languages
- universal dsl language xml
- domain specific

\section{Motivation}

% TODO: proc je uzitecne implementovat Neo4jCypher (zkusit pouziti DSL a MPS)



\section{Goals}

% TODO: zadani diplomky

The goals of this thesis are these:

\begin{itemize}
  \item Discuss MPS contribution in DSL languages
  \item Discuss MPS contribution in NoSQL domain
  \item Design and implement Neo4jCypher language and 
  		projectional editor and integrate it into MPS BaseLanguage (Java)
\end{itemize}
%%%%%%%%%%%%%%%%%%%%%%%%%%%%%%%%%%%%%
\chapter{Background}

\section{NoSQL databases}

\subsection{Neo4j graph database}

Neo4j is the most famous graph database system which is subset of NoSQL databases - databases which
offer more expresiveness than the classic SQL databases. The graph database is consisted of nodes
and relationships which gives us the power to query and update more complex expressions than in the
classic table database approach and gives us helpful graph functions like function for
searching shortest path and so. We chose this system because this system offers in addition to
classical java api approach also the DSL approach - Cypher DSL language which gives us the more
powerful tool to express what we want to process.

\subsection{Neo4j Cypher query language}
	
Cypher is a declarative domain specific query language for Neo4j graph database. It offers high
expressivity and efficient querying and updating of the graph in this database. It is designed both
for developers and for domain exprerts, so the syntax of the language is easy to understand.

\subsubsection{Syntax of Cypher}

To understand the principles and problems of design and implementation of Neo4jCypher (Cypher
implementation in MPS language workbench) we provide you the short overview of Cypher syntax.

Cypher query language is composed of these basic clauses:

\begin{itemize}
  \item \texttt{START} - specifies starting inputs
  \item \texttt{MATCH} - specifies pattern to match from starting nodes
  \item \texttt{WHERE} - next filtering of matched nodes
  \item \texttt{CREATE} - creates nodes, relationships
  \item \texttt{DELETE} - deletes nodes, relationships
  \item \texttt{SET} - updates nodes, relationships
  \item \texttt{RETURN} - returns nodes, relationships, etc.
\end{itemize}

Now we give some examples of queries with explanation what they do:

Read-Only Query:
\lstset{language=SQL}
\begin{lstlisting}
START
[MATCH]
[WHERE]
RETURN [ORDER BY] [SKIP] [LIMIT]
\end{lstlisting}

Write-Only Query:
\begin{lstlisting}
CREATE [UNIQUE]*
[SET|DELETE|FOREACH]*
[RETURN [ORDER BY] [SKIP] [LIMIT]]
\end{lstlisting}  


Read-Write Query:
\begin{lstlisting}
START
[MATCH]
[WHERE]
[CREATE [UNIQUE]]*
[SET|DELETE|FOREACH]*
[RETURN [ORDER BY] [SKIP] [LIMIT]]
\end{lstlisting}
	
%popsat jazyk Cypher (syntax, semantika, priklady, atd)

\section{Domain Specific Languages}

	\subsection{Neo4jCypher DSL}



\section{Language Workbenches}

	\subsection{How to define DSL in MPS}



\section{MPS Language Workbench}

	\subsection{How to Define DSL in MPS}

		\subsubsection{Structure DSL}
		\subsubsection{Editor DSL}
		\subsubsection{Constraints DSL}
		\subsubsection{Typesystem DSL}
		\subsubsection{Intensions and Other Parts}

	\subsection{MPS Pros and Cons}


%%%%%%%%%%%%%%%%%%%%%%%%%%%%%%%%%%%%%
\chapter{Design of Neo4jCypher}

\section{Problem analysis}

\section{Design decisions}

\subsection{Projection editor}

MPS Projection editor is a base approach of MPS to editing language. It differs from pure textual
editors in these points:

\begin{itemize}
  \item Unparsable editor notation - editor is composed of cells, every cell represents some node of
  the abstract syntax tree (AST)
  \item Possible switching among multiple notations using alternative editors, where every editor
  can show the AST in different way
  \item Combine notations from different authors and be extensible. It is possible to extend MPS
  languages, compose them and so on
\end{itemize}

\subsection{Text-like editor}

\subsection{Graphical extensions}


%%%%%%%%%%%%%%%%%%%%%%%%%%%%%%%%%%%%%
\chapter{Implementation details of Neo4jCypher language}

\section{Patterns}

\section{References}

\section{Integration into BaseLanguage}

MPS has a big advantage that we not only can design the new language and create the full ide
support for this language but also we can integrate this language into some existing language,
extend it and provide the way to execute our code with possibility to debug our code.

We chose these 4 steps, where each next step integrates the Neo4jCypher language more closely into
its BaseLanguage:

\begin{itemize}
  \item Design Neo4jCypher queries in Query Sheet editor - editor with full ide support:
  	\begin{itemize}
  	  \item code completition with correct reference visibility and full cypher function list
  	  \item intensions for query transformation, parts of query grouping, relationship
  	  type conversion, properties definition
  	  \item identifier uniquiness checking
  	  \item items type checking where necessary
  	\end{itemize} 
  \item Generation of MPCypher query into BaseLanguage string type - this approach provides us
  easy creation of query with full ide editor support. The next query manipulation depends on user's
  will. User can store this query or execute it using Neo4j Execution Engine from Neo4j java
  library.
  \item Using CypherExecutor object which adds us easy access to execution of Neo4jCypher query
  without need of implementation any further code from Neo4j java library. This approach tries to
  provide very easy way to query the database without need to know any king of Cypher java
  interface.
\end{itemize}



%%%%%%%%%%%%%%%%%%%%%%%%%%%%%%%%%%%%%
\chapter{Evaluation}

MPS platform is a robust tool for creating DSL languages


\section{Experience with MPS}

\section{MPS Contribution in DSL languages}

% Pros&Cons

\section{MPS Contribution in NoSQL Domain}

\section{Related work}
		
\section{Case Study}

% What to study in the future, next ideas of the work, improvements








\include{chap3_design}
\include{chap4_implementation}
\include{chap5_evaluation}

% Ukázka použití některých konstrukcí LateXu (odkomentujte, chcete-li)
% %%% Ukázka použití některých konstrukcí LaTeXu

\subsection{Ukázka \LaTeX{}u}
\label{ssec:ukazka}

This short subsection serves as an~example of basic \LaTeX{} constructs,
which can be useful for writing a~thesis.

Let us start with lists:

\begin{itemize}
\item The logo of Matfyz is displayed in figure~\ref{fig:mff}.
\item This is subsection~\ref{ssec:ukazka}.
\item Citing literature~\cite{lamport94}.
\end{itemize}

Different kinds of dashes:
red-black (short),
pages 16--22 (middle),
$45-44$ (minus),
and this is --- as you could have expected --- a~sentence-level dash,
which is the longest.
(Note that we have follwed \verb|a| by a~tilde instead of a~space
to avoid line breaks at that place.)

\newtheorem{theorem}{Theorem}
\newtheorem*{define}{Definition}	% Definice nečíslujeme, proto "*"

\begin{define}
A~{\sl Tree} is a connected graph with no cycles.
\end{define}

\begin{theorem}
This theorem is false.
\end{theorem}

\begin{proof}
False theorems do not have proofs.
\end{proof}

\begin{figure}
	\centering
	\includegraphics[width=30mm]{../img/logo.eps}
	\caption{Logo of MFF UK}
	\label{fig:mff}
\end{figure}


\chapter*{Conclusion}
\addcontentsline{toc}{chapter}{Conclusion}


%%% Seznam použité literatury
%%% Seznam použité literatury je zpracován podle platných standardů. Povinnou citační
%%% normou pro diplomovou práci je ISO 690. Jména časopisů lze uvádět zkráceně, ale jen
%%% v kodifikované podobě. Všechny použité zdroje a prameny musí být řádně citovány.

\def\bibname{Bibliography}
\begin{thebibliography}{99}
\addcontentsline{toc}{chapter}{\bibname}

%\bibitem{lamport94}
%  {\sc Lamport,} Leslie.
%  \emph{\LaTeX: A Document Preparation System}.
%  2. vydání.
%  Massachusetts: Addison Wesley, 1994.
%  ISBN 0-201-52983-1.

\bibitem{lamport94}
  {\sc Fowler,} Martin.
  \emph{Domain-Specific Languages}.
  Massachusetts: Addison Wesley, 2013.
  ISBN 0-321-71294-3.

\bibitem{lamport94}
  {\sc JetBrains}.
  \emph{MPS Documents and Live Demos}.
  \url{http://www.jetbrains.com/mps/documentation/index.html}.

\bibitem{lamport94}
  {\sc JetBrains}.
  \emph{MPS User's Guide}.
  \url{http://confluence.jetbrains.com/display/\\MPSD25/MPS+User's+Guide}.
  
\bibitem{lamport94}
  {\sc Voelter} Marcus.
  \emph{Language and IDE Modularization, Extension and Composition with MPS}.
  \url{http://voelter.de/data/pub/Voelter-GTTSE-MPS.pdf}.
  
\bibitem{lamport94}
  {\sc Neo Technology}.
  \emph{The Neo4j Manual v1.9.5}.
  \url{http://docs.neo4j.org/chunked/1.9.5/cypher-introduction.html}.
  
\end{thebibliography}


%%% Tabulky v diplomové práci, existují-li.
\chapwithtoc{List of Tables}

%%% Použité zkratky v diplomové práci, existují-li, včetně jejich vysvětlení.
\chapwithtoc{List of Abbreviations}

%%% Přílohy k diplomové práci, existují-li (různé dodatky jako výpisy programů,
%%% diagramy apod.). Každá příloha musí být alespoň jednou odkazována z vlastního
%%% textu práce. Přílohy se číslují.
\chapwithtoc{Attachments}

\openright
\end{document}
